\subsubsection{解析領域設定に伴う課題}
\label{sec:discussion_region_issues}

しかしながら、突発型 EEJ において Lunar age=18--21 付近では、
先行研究で報告されている傾向と比較して、
本研究では発生率がそれほど高くならない結果が得られた。

この結果の要因として、Dip 観測点と Offdip 観測点の経度差が
解析結果に影響を及ぼしている可能性が考えられる。
両観測点の経度差は約 $40^\circ$ であり、UTで2--3時間程度のずれが生じる。
その結果、磁気圏由来の変動が異なるタイミングで反映され、
EUEL の波形にずれが生じる可能性がある。
この影響により、未発達型と突発型の分類に用いた相関係数が低下し、
本来は未発達型として分類されるべきイベントが、
突発型として誤分類される可能性がある。


そこで本研究では、INTERMAGNET の TTB 観測点を用いることで、
観測点間の経度差を約 $10^\circ$ に縮小し、
磁気圏由来の変動差を低減することで、
EEJ の効果をより正確に捉えることを試みた。

次節では、このブラジル領域における解析結果に基づき、
未発達型 EEJ および突発型 EEJ の発生要因について考察する。